\chapter{Struktura programu}

\section{Dlaczego struktura jest wazna}
W malym pliku wszystko wydaje sie oczywiste. Problem zaczyna sie wtedy, gdy projekt
ma kilka stron, wspolne metadata, importy, funkcje pomocnicze i powtarzalne fragmenty.
Wlasnie wtedy porzadna struktura przestaje byc luksusem, a zaczyna byc koniecznoscia.

MoonChunk daje do tego kilka prostych narzedzi:

\begin{itemize}
  \item \mc{chunk} jako glowny blok organizacyjny,
  \item \mc{import} do ladowania rzeczy z innych plikow,
  \item \mc{@include} do wlaczania chunkow do chunkow,
  \item \mc{moon(...)} do wybierania punktu wejscia.
\end{itemize}

\section{Co moze znajdowac sie w pliku \mc{.mncnk}}
Plik MoonChunk zwykle sklada sie z dwoch czesci:

\begin{enumerate}
  \item opcjonalnych \mc{import}-ow na samej gorze,
  \item jednego lub wielu \mc{chunk}-ow oraz wywolan \mc{moon(...)}.
\end{enumerate}

Przyklad:

\begin{lstlisting}[style=moonchunk,caption={Typowy szkielet pliku MoonChunk}]
import { Shared } from "./shared.mncnk";
import * as Pages from "./pages.mncnk";

chunk "Main" {
  @include Shared;
  output: "./dist";
};

moon(Main);
\end{lstlisting}

Warto od razu zauwazyc dwie rzeczy:

\begin{itemize}
  \item importy stoja na samej gorze pliku,
  \item sam import nie uruchamia jeszcze generowania strony.
\end{itemize}

\section{Czym jest \mc{chunk}}
\mc{chunk} to podstawowa jednostka organizacyjna programu. Mozna go porownac do
modulu, ale z naciskiem na wykonanie. Chunk moze zawierac:

\begin{itemize}
  \item metadata,
  \item deklaracje zmiennych,
  \item funkcje,
  \item petle i warunki,
  \item strony \mc{page},
  \item wylaczenie innych chunkow przez \mc{@include}.
\end{itemize}

Przyklad:

\begin{lstlisting}[style=moonchunk,caption={Chunk zawierajacy kilka rodzajow elementow}]
chunk "Main" {
  output: "./dist";
  title: "Strona glowna";

  function greet(name: string): string {
    return "Czesc, " + name;
  }

  page "" {
    let userName = "Alicja";
    content {
      <h1>{title}</h1>
      <p>{greet(userName)}</p>
    };
  };
};
\end{lstlisting}

\section{Jedna nazwa, jedna odpowiedzialnosc}
Dobrze jest nadawac chunkom nazwy, ktore odpowiadaja ich roli. Przykladowo:

\begin{itemize}
  \item \mc{"Main"} -- glowny punkt wejscia,
  \item \mc{"Metadata"} -- wspolne metadata,
  \item \mc{"BlogPages"} -- generowanie podstron bloga,
  \item \mc{"ProductPages"} -- generowanie kart produktowych.
\end{itemize}

Taka konwencja bardzo pomaga przy importach i przy pracy zespolowej.

\section{Wywolanie \mc{moon(...)} jako punkt startowy}
Wazna cecha MoonChunk polega na tym, ze \mc{chunk} istnieje jako opis logiki,
ale dopiero \mc{moon(...)} powoduje jego wykonanie.

\begin{lstlisting}[style=moonchunk,caption={Jawny punkt wejscia}]
chunk "Main" {
  output: "./dist";
};

moon(Main);
\end{lstlisting}

Mozesz tez uruchomic kilka punktow wejscia:

\begin{lstlisting}[style=moonchunk,caption={Wiele wywolan moon(...)}]
chunk "Home" {
  output: "./dist";
  page "" {
    content {
      <h1>Home</h1>
    };
  };
};

chunk "About" {
  output: "./dist";
  page "about" {
    content {
      <h1>About</h1>
    };
  };
};

moon(Home);
moon(About);
\end{lstlisting}

Oznacza to, ze jeden plik moze wygenerowac wiele niezaleznych zestawow stron.

\section{Import a wykonanie}
To rozroznienie jest krytyczne. \important{Import nie znaczy wykonanie}. Import
sprawia tylko, ze jakis zasob staje sie dostepny w pliku. Dopiero pozniejsze
\mc{@include} lub \mc{moon(...)} powoduje realne uzycie danego fragmentu.

To zachowanie chroni przed nieprzewidywalnym uruchamianiem kodu z plikow pomocniczych.
Dzieki temu mozesz miec biblioteke chunkow i korzystac z niej tylko tam, gdzie chcesz.

\section{Eksportowanie chunkow}
Jesli chcesz importowac chunk z innego pliku, taki chunk powinien byc oznaczony jako
\mc{export}. Praktycznie wyglada to tak:

\begin{lstlisting}[style=moonchunk,caption={Chunk eksportowany do innych plikow}]
export chunk "Shared" {
  const siteName: string = "MoonChunk Docs";
};
\end{lstlisting}

Brak eksportu oznacza, ze dany chunk jest lokalny dla pliku i nie powinien byc
traktowany jak element publicznego interfejsu.

\section{Kolejnosc w pliku}
W MoonChunk kolejnosc ma znaczenie organizacyjne i semantyczne. Najbardziej bezpieczny
wzorzec wyglada tak:

\begin{enumerate}
  \item importy,
  \item deklaracje chunkow,
  \item wywolania \mc{moon(...)}.
\end{enumerate}

Wewnatrz samego \mc{chunk}-a dobrym zwyczajem jest trzymanie sie nastepujacego ukladu:

\begin{enumerate}
  \item \mc{@include},
  \item \mc{output} i metadata,
  \item \mc{env} z globalami,
  \item funkcje i stale,
  \item logika stron, petle, warunki, \mc{page}.
\end{enumerate}

Ten porzadek nie jest przypadkowy. Pozwala szybko rozpoznac, co jest wspolna konfiguracja,
co jest logika pomocnicza, a co odpowiada za faktyczne generowanie stron.

\section{Plik glowny i pliki pomocnicze}
W mniejszych projektach wszystko moze miescic sie w jednym pliku. W wiekszych duzo lepiej
sprawdza sie podzial na role, np.:

\begin{lstlisting}[style=basecode,caption={Przykladowy podzial projektu na pliki}]
content/
  site.mncnk
  shared.mncnk
  blog.mncnk
  pages/
    home.mncnk
    docs.mncnk
  data/
    posts.json
\end{lstlisting}

Wtedy:

\begin{itemize}
  \item \mc{site.mncnk} moze byc glownym plikiem startowym,
  \item \mc{shared.mncnk} moze przechowywac wspolne metadata i funkcje,
  \item pliki w \mc{pages/} moga zawierac chunki dla poszczegolnych obszarow witryny,
  \item katalog \mc{data/} moze trzymac zewnetrzne dane do ladowania przez \mc{data(...)}.
\end{itemize}

\section{Myslenie warstwowe}
Dobra struktura MoonChunk zwykle rozdziela projekt na trzy warstwy:

\begin{itemize}
  \item \important{warstwa konfiguracji} -- output, metadata, globalne wartosci,
  \item \important{warstwa logiki} -- funkcje, petle, warunki, obliczenia,
  \item \important{warstwa prezentacji} -- bloki \mc{content} i HTML.
\end{itemize}

Im wyrazniej rozdzielisz te warstwy, tym przyjemniej bedzie utrzymywac kod.

\section{Jak rozpoznac, ze struktura robi sie zla}
Oto kilka sygnalow ostrzegawczych:

\begin{itemize}
  \item jeden plik ma kilkaset linii i robi wszystko naraz,
  \item nie wiadomo, ktory \mc{chunk} jest punktem wejscia,
  \item metadata sa kopiowane do wielu miejsc,
  \item te same funkcje pojawiaja sie w kilku plikach,
  \item importy sa przypadkowe i trudno sledzic zaleznosci,
  \item chunki maja ogolne nazwy typu \mc{"Test"}, \mc{"Temp"}, \mc{"X"}.
\end{itemize}

Jesli widzisz takie objawy, warto zatrzymac sie na chwile i przebudowac uklad programu.
MoonChunk duzo zyskuje na czytelnej strukturze.
